
\textbf{LSM-based Key-Value Stores.} The evolution of technologies like 5G, AGI, and XR has led to an exponential increase in data generation and storage requirements. Many database systems have adopted Log-Structured Merge-tree (LSM) due to its efficiency in write operations. In LSM, data is initially written to a buffer; once the buffer is full, it is flushed to disk as an immutable Sorted String Table (SSTable) with sorted keys. LSM trees feature several levels, with each level being T times larger than the previous one. When a level reaches its capacity, multiple SSTable files are merged and sorted before being moved to the next level.

\textbf{Range Delete in LSM Trees.} Range deletes in LSM trees are handled by storing tombstones that specify the start and end keys, indicating the range of keys to be deleted in overlapping SSTable files. In systems like RocksDB, during compaction, these range delete tombstones can be used to purge obsolete entries across all compacted files.

\textbf{Problems.} 
A significant challenge in handling range deletes within SSTable files is the potential need to open and read specific files to ascertain the existence of an entry. This process can be slow. If it's possible to establish that a key does not exist in memory before accessing these files, we can expediently return and avoid the time-consuming task of reading files. This disparity in access speeds between disk and memory is a crucial factor that can decelerate query response times. By identifying non-existent keys earlier, unnecessary disk reads can be circumvented, thereby enhancing the overall efficiency of the system.


\subsection{Motivation}
By storing range delete information in the memory,  an RDF offers a streamlined approach for handling queries. It actively filters out queries aimed at keys already marked for deletion, thereby preventing the unnecessary activation of disk-based operations. This approach is particularly beneficial in scenarios where the frequency of range deletions is high, and point queries often intersect with these ranges.

Moreover, this memory-based RDF implementation opens up new possibilities for optimizing the LSM tree's operational dynamics. By allowing for rapid determination of the non-existence of keys, the RDF effectively narrows down the search space for point queries, thus speeding up the overall query response time. This optimization is especially significant in systems with large datasets and frequent read operations, where disk reads constitute a substantial operational bottleneck.

In summary, it not only reduces the number of disk I/O operations during point queries but also enhances the efficiency of the entire system by enabling quicker and more effective data searching processes.

\subsection{Problem Statement}
The search for a key, given the presence of range delete data in SSTable files, may be inefficient on LSM storage. Point queries initiate their searching in memory and then progress to disk if necessary. The organizational structure of these files spans multiple levels, and point queries systematically move from the lower to higher levels in this hierarchy. When a query is made, a relevant SSTable file on each level is targeted when its range, spanning from its minimum key to its maximum key, encompasses the queried key. The process involves reading range delete information from the file's metadata to ascertain existence of the key.  However, this approach can lead to inefficiencies, particularly when there's repetitive reading of range delete data or when multiple files across different levels are opened to retrieve this information and search for the key, resulting in unnecessary I/O operations.

A more efficient approach would be to store range delete information directly in memory. This method could sidestep the need for opening certain files , thereby reducing I/O operations. Furthermore, considering range delete and point query operations coming in order of their issued sequence, associating sequence numbers with range delete tuples, the pair of start and end keys, becomes redundant because data on the lower level is always newer than the ones on the higher level.  Eliminating the use of sequence numbers in Range Delete Filters (RDFs) allows for the merging of overlapping ranges within the same level, thereby further enhancing system efficiency and reducing the operational overhead.



\subsection{Contributions}
This study aims to:
(a) Discuss and analyze the impact of different RDF implementations (PLRDF, split-PLRDF, toplevel-RDF, Skyline-RDF) on system performance and their memory footprints.
(b) Eliminate the sequence number associated with each range to enable range merging in RDFs to decrease its memory footprint. 